\section{引言}
\label{sec:introduction}

\subsection{研究背景及意义}

在“双碳”目标与新型电力系统建设持续推进的背景下，我国电源结构正发生显著变化。国家能源局发布的2025年全国电力统计数据显示，截至2025年底，全国累计发电装机容量达到38.9亿kW，其中太阳能发电装机容量达到12.0亿kW，风电装机容量达到6.4亿kW\cite{NEA2026PowerStatistics}；国家统计局进一步指出，2025年可再生能源装机总量达到23.4亿kW，火电装机容量约15.4亿kW，占全国发电装机容量的比重首次降至40\%以下\cite{NBS2026EnergyTransition}。随着风电、光伏等波动性新能源大规模接入，电力系统对可调节电源的快速响应、深度调峰和宽负荷运行能力提出了更高要求。与此同时，煤电在我国电力安全保供和系统调节中仍承担重要的支撑作用，其功能定位正在由传统主体电源逐步向基础保障性、系统调节性电源转变。国家发展改革委、国家能源局发布的《新一代煤电升级专项行动实施方案（2025--2027年）》明确提出推进煤电向清洁低碳、安全可靠、高效调节和智能运行方向升级，并进一步强化快速变负荷、深度调峰和启停调峰能力\cite{NDRC2025CoalPowerUpgrade}。因此，在高比例新能源电力系统中，如何保证燃煤机组在频繁变工况运行条件下仍具有较高的安全性、可靠性和运行经济性，已成为煤电智能化转型过程中需要重点解决的问题。

燃煤机组灵活运行能力的提升并非单纯依赖锅炉侧调节，其本质上涉及锅炉、汽轮机及回热系统等多环节之间的动态能量协调。已有研究表明，随着可再生能源渗透率提升，传统燃煤机组需要在更低最小负荷、更高爬坡速率以及更频繁负荷循环条件下运行，机组偏离设计工况的程度显著增加\cite{GonzalezSalazar2018Flexibility}。汽轮机作为燃煤机组实现蒸汽热能向机械能转换的核心设备，是负荷调节过程中直接响应蒸汽流量与热力边界变化的关键环节。Yu等针对330~MW亚临界机组和1000~MW超超临界机组的研究指出，新能源并网背景下火电机组非设计工况运行趋于频繁，而汽轮机调节级的流量特性、效率特性及出口参数直接关系到机组灵活运行过程中的效率与安全\cite{Yu2020HybridDigitalTwin}。由此可见，汽轮机系统运行状态能否被准确表征，是燃煤机组实现宽负荷安全运行和灵活调节的重要基础。

大型燃煤机组汽轮机并非单一设备，而是由多级汽缸、再热系统、抽汽回热系统、凝汽器及其相关管路共同组成的强耦合热力系统。以二次再热汽轮机为例，主蒸汽首先进入超高压缸做功，排汽经一次再热后进入高压缸，再经二次再热进入中压缸和低压缸，并通过多级抽汽向回热系统提供热源，最终由凝汽器建立低压缸排汽边界。不同设备之间通过蒸汽流动、压力传递、能量交换以及控制调节形成显著的上下游耦合关系。Zhang等基于600~MW燃煤机组构建了汽轮机全工况模型，将汽轮机级组、再热器、抽汽管路、加热器及凝汽器等划分为多个耦合子模型，并利用实际运行数据验证了模型对不同负荷运行参数以及典型设备异常状态的表征能力\cite{Zhang2018AllCondition}。该研究说明，汽轮机系统状态并不能由单一测点或单一设备独立刻画，而应结合运行边界和设备间热力关联对系统状态进行整体分析。

% [图1占位] 燃煤机组汽轮机系统及主要运行边界示意图
% 建议内容：主蒸汽 -> VHP -> 一次再热 -> HP -> 二次再热 -> IP -> LP -> 凝汽器；
% 同时标示抽汽回热支路以及主蒸汽参数、再热蒸汽参数、机组负荷、凝汽器背压等外部运行边界。

频繁升降负荷和偏离设计工况运行还会增加汽轮机关键部件状态变化与性能退化识别的难度。一方面，负荷、主再热蒸汽参数、抽汽状态以及凝汽器背压等正常运行边界变化本身即可引起汽缸压力、温度、流量及效率参数的大幅波动，使固定阈值或设计工况基准难以直接用于长期在线评价；另一方面，快速负荷变化和频繁启停会提高关键部件的热机械载荷水平。Banaszkiewicz针对汽轮机转子开展的在线热应力监测研究表明，高灵活性运行条件下启停过程和负荷变化会导致关键部件承受显著热应力与低周疲劳载荷，因此有必要利用在线测量和数学模型持续监测关键状态\cite{Banaszkiewicz2016ThermalStress}。此外，汽轮机在长期运行过程中还可能出现叶片表面粗糙度增加、固体颗粒侵蚀、汽封磨损及叶片沉积等通流部件退化。Kubiak等通过实际汽轮机检修案例表明，这类部件退化能够通过关键压力及通流性能参数的变化在运行阶段被识别，并在后续检修中得到验证\cite{Kubiak2002TurbineDegradation}。因此，面向长期在线运行建立能够区分正常工况变化与真实设备状态偏离的监测方法，对汽轮机早期异常识别和检修决策具有重要工程意义。

传统汽轮机性能监测通常建立在热平衡、级组效率和流量特性等机理计算基础上，具有较好的物理解释性，但其在线应用精度容易受到测量不确定性、结构参数缺失以及非设计工况模型误差影响。Jiang等在燃煤机组汽轮机在线性能监测研究中指出，蒸汽循环热平衡分析高度依赖流量等关键测量的准确性，并利用数据协调方法降低测量不确定性对热耗率和性能监测结果的影响\cite{Jiang2014DataReconciliation}。与此同时，随着分布式控制系统（Distributed Control System, DCS）和厂级监控信息系统长期积累高频运行数据，利用数据驱动模型学习运行边界与设备状态之间的正常映射关系成为可能。Dettori等分别采用热力学混合模型和神经网络模型开展汽轮机在线监测建模，并利用高压汽轮机现场数据验证了模型预测能力\cite{Dettori2017MonitoringModels}；Yu等进一步将机理信息与运行数据融合用于全工况汽轮机调节级数字孪生，实现了不同运行范围内出口压力和温度等关键参数的高精度预测\cite{Yu2020HybridDigitalTwin}。这些研究表明，基于运行数据建立正常状态参考模型，可以为复杂变工况下汽轮机状态监测提供不同于固定设计基准的动态健康参照。

从故障诊断角度看，汽轮机局部异常往往不会局限于故障源所在的单一参数，而可能沿蒸汽通流、抽汽回热和压力边界关系向上下游设备传播，最终表现为多个热工参数的协同偏离。因此，仅依赖单测点阈值或将全部测点视为相互独立的特征变量，难以完整描述复杂故障的传播过程。面向燃煤机组汽轮机系统开展状态监测与故障诊断研究，一方面需要建立适应宽负荷运行条件的健康状态模型，在给定运行边界下获得关键参数的正常参考值，从而削弱正常工况迁移对异常识别的干扰；另一方面还需要充分利用汽轮机设备拓扑、介质流向及故障机理等先验知识，描述异常参数之间具有物理意义的结构关联，为后续故障类型识别和传播路径分析提供依据。由此构建从“正常状态表征—异常偏离提取—系统关联诊断”逐层递进的技术体系，不仅能够提升燃煤机组汽轮机在复杂运行条件下的状态感知能力，也可为故障早期预警、故障定位和状态检修提供数据与知识支撑，具有重要的理论研究意义和工程应用价值。

\subsection{燃煤机组汽轮机系统状态监测及故障诊断研究现状}

\subsubsection{基于机理驱动的状态监测及故障诊断方法}

基于机理驱动的状态监测及故障诊断方法以设备结构、热力过程和守恒关系为基础，通过建立能够表征健康状态的数学模型，将实际测量结果与模型计算结果进行比较，并依据效率、流量、压力、温度或残差等指标判断设备状态。汽轮机系统具有明确的质量流动和能量转换过程，因此质量守恒、能量守恒、蒸汽热力性质、级组效率以及通流能力等长期以来都是其性能评价和故障诊断的重要理论基础。与完全依赖统计相关性的方法相比，机理驱动方法能够将状态变化对应到具体热力过程或部件性能变化，其诊断结果通常具有较明确的工程物理含义。

早期汽轮机热力诊断研究主要通过建立正常参考状态并比较实际性能偏差来识别部件异常。Zaleta-Aguilar等提出面向动力系统部件的热力表征方法，利用焓变、熵变及质量流量比等参数构造参考性能状态（Reference Performance State, RPS），并在汽轮机、凝汽器、换热器和泵等对象上分析部件内部退化及外部诱发异常对参考状态的偏离\cite{Zaleta2004ThermoCharacterization}。该类方法的核心思想是先依据设计信息、热力过程与仿真计算确定无故障状态，再利用实际运行参数与健康基准之间的偏差进行状态评价。Kubiak等进一步针对叶片粗糙度、喷嘴或叶片侵蚀、汽封间隙变化和沉积等汽轮机通流部件退化开展案例研究，通过关键压力和简化流量关系识别部件性能变化，并利用停机检修结果对诊断结论进行了验证\cite{Kubiak2002TurbineDegradation}。这说明以通流能力、压力分布和级组性能变化为核心的热力参数能够较直接地反映汽轮机内部通流状态。

随着机组运行工况不断扩展，单一设计点下的性能基准难以覆盖宽负荷条件，研究重点逐渐由静态热平衡诊断向全工况动态模型与在线状态估计发展。Zhang等依据汽轮机各级组及回热系统的物理过程建立了600~MW机组全工况仿真模型，将汽轮机、再热器、抽汽管路、给水加热器和凝汽器等多个模块耦合起来，实现不同运行负荷下系统参数的动态计算\cite{Zhang2018AllCondition}。Yu等则围绕汽轮机调节级构建机理与数据融合的数字孪生模型，根据阀门流量关系和级组热力特性计算调节级状态，并通过运行数据辨识和修正模型参数，使模型能够用于不同工况下出口压力、温度等性能参数的在线监测\cite{Yu2020HybridDigitalTwin}。对于实际机组测量误差问题，Jiang等将非线性数据协调引入汽轮机在线性能监测，在质量与能量平衡等约束下对运行测量数据进行校正，以提高热耗率及相关性能参数计算的可信度\cite{Jiang2014DataReconciliation}。上述工作表明，利用机理约束建立随运行边界变化而变化的正常状态模型，是提高宽工况汽轮机性能监测准确性的有效途径。

除直接计算热力性能指标外，基于状态观测器和残差分析的模型诊断方法也被用于锅炉—汽轮机耦合系统。Madrigal-Espinosa等基于实际电厂低、中、高负荷数据建立准线性参数时变（quasi-linear parameter varying, quasi-LPV）锅炉—汽轮机模型，并设计Luenberger型观测器产生模型残差，通过自适应阈值完成关键传感器故障检测与隔离\cite{Madrigal2017BoilerTurbineFDI}。该研究说明，模型预测与实际观测之间的残差不仅可以反映系统状态偏离，还可以通过自适应基准减少负荷变化造成的误报警。Banaszkiewicz则从热机械安全角度利用实时温度和机理模型在线估计汽轮机转子热应力，为启停及变负荷过程中的寿命消耗监测和运行控制提供依据\cite{Banaszkiewicz2016ThermalStress}。因此，机理模型既可以直接形成性能指标，也可以作为生成健康参考值和状态残差的基础。

总体而言，机理驱动方法能够充分利用汽轮机质量守恒、能量守恒和通流特性等先验规律，在故障机理解释、状态参数推演和诊断结果可信度方面具有明显优势。然而，大型汽轮机系统包含多级汽缸、再热器、抽汽回热设备及凝汽器等多个耦合部件，精细模型往往需要较完整的设备结构、设计工况、阀门特性、级组参数和蒸汽物性信息；实际运行中的测量误差、设备老化、检修改造及长期性能漂移又会造成模型参数与真实对象逐渐偏离。虽然数据协调、参数辨识及混合建模能够缓解上述问题\cite{Jiang2014DataReconciliation,Yu2020HybridDigitalTwin}，但对于需要长期在线部署的大规模机组而言，逐设备建立并持续修正高精度机理模型仍具有较高的专业知识和维护成本。因此，在保留必要物理边界和设备拓扑信息的同时，利用长期运行数据自主学习复杂非线性状态关系，成为汽轮机状态监测与故障诊断的重要研究方向。

\subsubsection{基于数据驱动的状态监测及故障诊断方法}

随着DCS、汽轮机监视仪表系统（Turbine Supervisory Instrumentation, TSI）及厂级信息系统的持续建设，燃煤机组积累了大量包含负荷、压力、温度、流量、阀位以及设备状态信息的历史运行数据。基于数据驱动的方法无需完整求解汽轮机内部复杂热力过程，而是利用历史样本建立输入运行条件与目标状态参数之间的统计或非线性映射，或者直接学习健康数据在高维特征空间中的正常分布。按照模型任务不同，相关研究主要可以分为正常状态建模与参数预测、无监督异常检测以及有监督故障分类等方向。

正常状态建模的基本思想是在健康运行数据上学习设备随工况变化的正常响应关系，以动态参考值替代固定设计值。Cao等针对凝汽式汽轮机末级组湿蒸汽区排汽焓难以准确获取的问题，分析影响末级组相对内效率的运行因素，并采用组合BP神经网络建立正常相对内效率计算模型，其宽工况测试结果与理论计算结果具有较好一致性，为末级组状态监测提供了数据驱动正常值基准\cite{Cao2016LastStageNormalValue}。Dettori等比较了热力学混合模型和神经网络模型在汽轮机在线监测中的应用，利用高压汽轮机现场运行数据预测关键状态参数，说明神经网络可以在缺少部分直接测量量时学习复杂设备的正常运行映射\cite{Dettori2017MonitoringModels}。这些工作表明，将运行工况作为模型条件并预测设备正常状态，是削弱非设计工况影响的一种可行思路。

汽轮机运行数据具有明显的时间相关性，频繁升降负荷时当前状态不仅与瞬时输入有关，还受到前序运行过程和热惯性的影响。因此，循环神经网络等时序模型逐渐被用于汽轮机动态状态监测。Li等针对传统稳态评价方法在动态运行过程中的适应性不足问题，将机理分析与数据特征选择相结合，并利用长短期记忆网络（Long Short-Term Memory, LSTM）建立级组性能退化预警模型；实验结果显示，考虑时间序列信息有助于提升动态工况下性能退化识别能力\cite{Li2021LSTMStageWarning}。近年来，针对多工况状态迁移问题，Chen等构建了基于多高斯过程自适应迁移学习的汽轮机在线无监督状态监测框架，将环境变化、功率需求变化及人工调整引起的不同运行模式显式纳入监测过程，并通过两级判据区分正常模式迁移与异常状态；该方法在实际汽轮机数据上完成验证，且无需故障标签\cite{Chen2023AdaptiveTransfer}。这类研究说明，宽工况状态监测的关键不仅是提高单一工况下的预测精度，还需要避免将正常工况迁移错误判断为设备异常。

考虑到电站重大设备故障发生频率低且出现异常后通常会及时停机或处置，真实汽轮机故障样本往往远少于正常样本，无监督或仅依赖健康数据的异常检测方法受到广泛关注。Ko等采用集成降噪自编码器建立汽轮机正常状态模型，并根据模型输出与残差联合分布构造动态阈值，以降低正常状态波动造成的误报警，同时利用残差敏感度识别与异常相关的监测信号\cite{Ko2022DynamicThresholdAE}。Zhai等利用真实660~MW机组运行数据构建基于卷积自编码器的汽轮机故障预警模型，通过模型预测输出与实际运行数据之间的偏差识别异常，并结合混合信息增益筛选多维传感器特征\cite{Zhai2024CAEWarning}。Xu和Zhang进一步采用增强LSTM变分自编码器对高维汽轮机时序数据进行无监督表示学习，将潜变量和重构误差融合为异常特征，并在实际工业汽轮机数据上进行验证\cite{Xu2025LSTMVAE}。从这些研究可以看出，基于正常数据学习健康行为并利用预测或重构偏差表征异常，能够在缺少完整故障标签的情况下充分利用电厂长期积累的健康运行样本。

在明确故障类型的诊断任务中，支持向量机、神经网络、模糊推理及其组合模型也被广泛应用。Salahshoor等面向工业汽轮机构建SVM与自适应神经模糊推理系统（Adaptive Neuro-Fuzzy Inference System, ANFIS）融合的故障检测与诊断方法，通过有序加权平均实现多分类器信息融合，并在包含高压、中压和低压部分的工业汽轮机仿真对象上验证了多类早期故障识别能力\cite{Salahshoor2010SVMANFIS}。此类监督分类方法能够直接建立多变量运行特征与故障类别之间的映射，在训练故障模式与实际测试模式一致时通常具有较高分类效率，但其性能较大程度取决于故障样本覆盖范围、类别均衡程度以及训练和实际运行工况之间的数据分布一致性。

传统神经网络和机器学习模型通常将多测点数据组织为规则向量或矩阵进行特征学习，而汽轮机的汽缸、再热器、抽汽系统和凝汽器之间具有明确的物理连接关系，故障影响也可能沿蒸汽流动和设备连接关系逐级传播。为显式利用复杂系统中的结构信息，近年来图神经网络（Graph Neural Network, GNN）逐渐被引入工业故障诊断。Zhang和Hu针对汽轮机系统首先依据系统机理建立过程图，并进一步构建包含潜在故障节点和传播关系的故障图，采用GNN学习故障节点表示并通过链路预测推断汽轮机泄漏故障传播路径，验证了以图结构表达汽轮机设备关系和故障传播的可行性\cite{ZhangHu2021FaultPropagationGNN}。在通用工业过程领域，Jia等根据工艺流程图建立变量拓扑图，将自注意力、图卷积和图池化结合用于故障诊断，使模型特征传播过程受到实际流程拓扑约束\cite{Jia2023TopologyGuidedGraph}；Han等则将过程知识模型与GNN结合，通过知识约束减少因纯数据因果发现产生的冗余连接，并用于复杂工业过程的故障根因分析\cite{Han2022KnowledgeEnhancedGNN}。这些研究表明，将设备拓扑、工艺流程及领域知识转化为图结构，可以使数据驱动模型在提取统计特征之外进一步利用具有物理意义的参数关联关系，为复杂热力系统的故障传播建模和诊断解释提供新的技术路径。

综上，数据驱动方法在非线性拟合、时序特征学习和复杂高维数据处理方面具有明显优势，并已经从传统监督分类逐渐发展到健康状态建模、无监督异常检测和知识增强图学习等多种技术路线。然而，面向燃煤机组汽轮机系统的长期在线应用仍存在若干不足：一是汽轮机参数随负荷和热力边界显著变化，若直接使用原始参数进行异常检测或故障分类，正常工况迁移与真实故障造成的状态偏离可能产生特征混叠\cite{Li2021LSTMStageWarning,Chen2023AdaptiveTransfer}；二是高质量故障样本获取困难，仅依赖监督学习难以充分利用长期积累的大量健康运行数据\cite{Ko2022DynamicThresholdAE,Chen2023AdaptiveTransfer}；三是多数基于规则向量的数据驱动模型未显式表达汽轮机设备之间的蒸汽流动、能量传递和故障传播关系，而现有图学习研究虽然证明了拓扑知识在工业故障诊断中的价值\cite{ZhangHu2021FaultPropagationGNN,Jia2023TopologyGuidedGraph,Han2022KnowledgeEnhancedGNN}，但如何将汽轮机健康状态模型得到的动态异常信息与系统级领域知识进行统一融合，仍具有进一步研究和工程应用价值。

\subsection{现有研究与应用存在的问题}

根据上述研究现状，现有燃煤机组汽轮机状态监测及故障诊断方法已在热力性能计算、正常行为建模、异常检测及智能分类等方面取得较多成果，但在复杂变工况和系统级故障诊断场景下仍存在以下问题。

（1）\textbf{正常工况变化与设备异常在原始参数空间中存在混叠。} 汽轮机压力、温度、流量及抽汽状态不仅受设备健康程度影响，还随负荷、主再热蒸汽参数、凝汽器背压及运行策略改变而发生显著变化。已有研究已经从动态性能建模和多运行模式监测角度指出，传统稳态或单模式模型在频繁负荷变化条件下容易降低监测灵敏度，正常模式切换也可能造成误报警\cite{Li2021LSTMStageWarning,Chen2023AdaptiveTransfer}。因此，仅依据实际测量值的绝对变化识别故障难以稳定适用于宽负荷运行场景，需要首先建立能够表征给定运行边界下健康响应的正常状态模型，并利用实际状态相对健康参考的偏离量提取异常特征。

（2）\textbf{纯数据驱动诊断模型难以显式利用汽轮机系统的物理拓扑和故障传播知识。} 大型汽轮机由超高压缸、高压缸、中压缸、低压缸、再热系统、抽汽回热系统和凝汽器等多个设备构成，参数之间通过蒸汽流动、压力传递和能量交换形成明确结构关系。当局部故障发生时，异常特征可能沿设备连接逐层传播。现有向量化机器学习模型主要依赖样本中的统计相关性，其内部特征交互不必然与真实热力拓扑一致。汽轮机故障传播GNN及流程拓扑引导图学习研究已经证明，显式引入系统结构能够为故障特征聚合和传播路径推断提供有效约束\cite{ZhangHu2021FaultPropagationGNN,Jia2023TopologyGuidedGraph}。因此，有必要进一步将汽轮机设备、参数、异常现象和故障机理等运行知识结构化表达，并作为诊断模型的拓扑先验。

（3）\textbf{真实高质量故障样本稀缺，监督式诊断模型训练受到限制。} 汽轮机属于高可靠性重大旋转设备，异常状态通常会被及时处置，重大故障发生频率较低，导致现场长期数据中健康样本数量远高于有效故障样本。Ko等和Chen等均指出，汽轮机真实故障数据难以充分获取，因而采用仅依赖正常数据或不依赖故障标签的无监督监测策略\cite{Ko2022DynamicThresholdAE,Chen2023AdaptiveTransfer}。若直接利用少量有标签故障样本从零开始训练高容量深度分类模型，容易受到样本数量和类别分布影响。因此，需要在故障分类之前充分利用大量无标签健康运行数据学习汽轮机参数之间的正常关联结构，再利用有限故障样本完成判别边界训练。

（4）\textbf{健康状态监测与系统级故障诊断之间缺少统一的信息衔接。} 现有研究中，一类方法侧重于建立正常行为模型并依据预测误差或重构误差检测异常\cite{Dettori2017MonitoringModels,Ko2022DynamicThresholdAE,Zhai2024CAEWarning}，另一类方法则直接基于故障样本建立分类器或利用图结构推断故障关系\cite{Salahshoor2010SVMANFIS,ZhangHu2021FaultPropagationGNN}。两类研究分别解决“运行状态是否偏离健康状态”和“异常属于何种故障或如何传播”的问题，但在实际汽轮机系统中二者具有天然的前后联系。若能够将健康状态模型生成的参数偏差作为动态故障表征，并进一步在汽轮机知识拓扑中进行关联特征聚合，即可形成从正常状态基准、异常特征提取到故障类别识别的连续诊断链路。

基于上述问题，本文面向燃煤机组汽轮机系统构建“健康状态建模—动态残差表征—知识图谱约束—图卷积故障诊断”的统一技术框架。首先利用正常运行数据建立适应运行边界变化的设备级健康模型，将实际测量值与健康模型预测值之间的偏差作为状态异常表征；随后利用运行生产资料建立汽轮机故障知识图谱，将设备拓扑和参数关联关系转化为图卷积网络可计算的邻接结构；最后在图拓扑约束下聚合多参数动态残差，实现故障类别识别及关联特征分析。

\subsection{论文主要研究内容}

本文以某燃煤机组汽轮机系统为研究对象，围绕复杂变工况下正常状态表征不足、系统拓扑知识利用不充分及故障样本有限等问题，研究基于健康状态模型与知识图谱图卷积网络的汽轮机故障诊断方法。全文总体技术路线如图~\ref{fig:overall_framework}所示，主要研究内容如下。

（1）\textbf{构建基于GRU的燃煤机组汽轮机系统健康状态监测方法。} 根据汽轮机热力系统的介质流动关系和设备边界，将运行操作变量与上游/外部热力边界作为模型输入，以反映设备正常状态的关键压力、温度等参数作为模型输出，建立设备级健康状态预测模型。考虑汽轮机运行参数具有动态迟延和时间相关特性，引入门控循环单元（Gated Recurrent Unit, GRU）学习历史时间窗口内的多变量动态映射，并设置ANN、TCN及iTransformer等模型进行对比。以超高压缸为典型对象开展工程算例，通过实际运行数据验证健康模型对复杂运行边界下正常参数变化的表征能力，并进一步由实际观测值与模型预测值构造标准化动态残差，为后续故障诊断提供状态特征。

（2）\textbf{构建面向汽轮机系统故障诊断的运行知识图谱。} 面向汽轮机本体及相关热力系统，综合设备说明、运行规程、系统拓扑、故障案例及相关运行生产资料，对设备、监测参数、异常现象、故障类型和处置等知识实体进行统一定义；通过文本知识抽取、关系校验和实体标准化等过程形成结构化三元组，建立覆盖超高压缸、高压缸、中压缸、低压缸、再热及抽汽回热等关键环节的汽轮机故障知识图谱。进一步从知识图谱中提取与状态监测参数对应的关联关系，构建可用于图卷积计算的参数邻接矩阵。

（3）\textbf{提出融合健康状态残差与知识图谱的KG-GCN故障诊断方法。} 将健康模型产生的连续标准化残差按照滑动时间窗口映射为知识图谱中参数节点的动态特征，并以知识图谱得到的参数邻接矩阵作为图卷积网络的结构输入，使不同参数的异常信息沿具有实际工程意义的关联路径进行特征传播与聚合。针对故障标签有限的问题，首先利用大量无标签运行残差窗口对图卷积编码器进行掩码重构预训练，使模型学习汽轮机参数在知识拓扑约束下的关联特征，再利用有限有标签故障样本训练故障分类层，实现汽轮机系统故障识别。

（4）\textbf{开展典型对象故障诊断实验并验证方法有效性。} 以超高压缸及其关联热力参数为典型对象构建故障诊断算例，采用准确率、精确率、召回率及Macro-F1等指标评价诊断性能，并与传统数据驱动诊断模型进行对比。同时通过去除健康残差表征或知识拓扑约束等消融实验，分别分析健康状态建模与知识图谱对故障诊断性能的贡献，并结合异常参数及图节点关联关系分析典型故障的特征传播过程。

\begin{figure}[htbp]
    \centering
    % \includegraphics[width=0.96\textwidth]{figures/overall_framework.pdf}
    \fbox{\parbox[c][5.0cm][c]{0.92\textwidth}{\centering
    图占位：基于健康状态建模与KG-GCN的燃煤机组汽轮机系统故障诊断总体框架\\[0.5em]
    DCS健康运行数据 $\rightarrow$ GRU健康模型 $\rightarrow$ 标准化动态残差\\
    运行规程/说明书/故障资料 $\rightarrow$ 知识抽取 $\rightarrow$ 汽轮机故障知识图谱\\
    动态残差窗口 + 参数邻接矩阵 $\rightarrow$ KG-GCN $\rightarrow$ 故障类别与关联特征分析}}
    \caption{基于健康状态建模与KG-GCN的燃煤机组汽轮机系统故障诊断总体框架}
    \label{fig:overall_framework}
\end{figure}
